\documentclass[utf8, 12pt]{bmstu-report}

\usepackage{listings}
\usepackage{xcolor}
\usepackage{booktabs}

% Настройка подсветки кода SQL
\lstset{
	language=SQL,
	basicstyle=\small\ttfamily,
	keywordstyle=\color{blue}\bfseries,
	commentstyle=\color{gray},
	stringstyle=\color{teal},
	numbers=left,
	numberstyle=\tiny\color{gray},
	stepnumber=1,
	numbersep=5pt,
	frame=single,
	breaklines=true,
	captionpos=b,
	keepspaces=true
}

\author{И. И. Иванов}
\group{ИУ7-51B}
\title{Лабораторная работа № 3}
\subtitle{Запросы PostgreSQL}
\teacher{Ю. А. Григорьев}
\department{ИУ7 "Программное обеспечение ЭВМ и информационные технологии"}

\begin{document}
	
	\maketitle
	
	\tableofcontents
	\newpage
	
	\section{Цель работы}
	Целью лабораторной работы является формирование навыков разработки сложных SQL-запросов SELECT с использованием соединений таблиц, агрегатных функций, группировки, оконных функций и подзапросов в СУБД PostgreSQL.
	
	\section{Выполнение контрольных заданий}
	
	\subsection{Задание 1. Правила обработки пробелов в операторе LIKE}
	Было произведено сравнение следующих запросов:
	\begin{lstlisting}
		SELECT count(*) FROM tickets;
		SELECT count(*) FROM tickets WHERE passenger_name LIKE '% %';
		SELECT count(*) FROM tickets WHERE passenger_name LIKE '% % %';
		SELECT count(*) FROM tickets WHERE passenger_name LIKE '% %%';
	\end{lstlisting}
	
	\textbf{Пояснение:} 
	Символ \texttt{\%} в операторе \texttt{LIKE} обозначает любую последовательность символов, включая пустую. В связи с этим шаблон \texttt{'\% \%\%'} полностью эквивалентен \texttt{'\% \%'}. Шаблон \texttt{'\% \% \%'} ищет строки, содержащие минимум два пробела.
	
	\subsection{Задание 2. Выбор пассажиров с фамилией из 5 букв}
	Запрос для поиска пассажиров, у которых фамилия состоит из 5 букв:
	\begin{lstlisting}
		SELECT passenger_name
		FROM tickets
		WHERE passenger_name LIKE '% _____';
	\end{lstlisting}
	
	\subsection{Задание 6. Маршруты на самолетах Boeing}
	Запрос для поиска маршрутов, обслуживаемых самолетами Boeing, с выводом наименования модели:
	\begin{lstlisting}
		SELECT DISTINCT
		r.departure_city,
		r.arrival_city,
		a.model
		FROM routes r
		JOIN aircrafts a ON r.aircraft_code = a.aircraft_code
		WHERE a.model LIKE 'Boeing%';
	\end{lstlisting}
	
	\subsection{Задание 7. Города, между которыми летают Boeing 777-300}
	Запрос для вывода уникальных пар городов:
	\begin{lstlisting}
		SELECT DISTINCT
		LEAST(r.departure_city, r.arrival_city) AS city1,
		GREATEST(r.departure_city, r.arrival_city) AS city2
		FROM routes r
		JOIN aircrafts a ON r.aircraft_code = a.aircraft_code
		WHERE a.model = 'Boeing 777-300';
	\end{lstlisting}
	
	\subsection{Задание 8. Количество рейсов из Москвы в Санкт-Петербург}
	\begin{lstlisting}
		SELECT departure_city, arrival_city, count(*) AS count
		FROM routes
		WHERE departure_city = 'Москва' AND arrival_city = 'Санкт-Петербург'
		GROUP BY departure_city, arrival_city;
	\end{lstlisting}
	
	\subsection{Задание 10. Количество направлений из каждого города}
	\begin{lstlisting}
		SELECT departure_city, count(DISTINCT arrival_city) AS num_directions
		FROM routes
		GROUP BY departure_city
		ORDER BY num_directions DESC;
	\end{lstlisting}
	
	\subsection{Задание 11. Топ-5 городов по числу ежедневных рейсов из Москвы}
	\begin{lstlisting}
		SELECT arrival_city, count(*) AS daily_routes_count
		FROM routes
		WHERE departure_city = 'Москва' 
		AND array_length(days_of_week, 1) = 7
		GROUP BY arrival_city
		ORDER BY daily_routes_count DESC
		LIMIT 5;
	\end{lstlisting}
	
	\subsection{Задание 12*. Число рейсов из Москвы по дням недели}
	\begin{lstlisting}
		SELECT 
		unnested_day AS day_of_week,
		count(*) AS flights_count
		FROM routes,
		unnest(days_of_week) AS unnested_day
		WHERE departure_city = 'Москва'
		GROUP BY unnested_day
		ORDER BY unnested_day;
	\end{lstlisting}
	
	\subsection{Задание 13. Максимальные и минимальные цены билетов}
	\begin{lstlisting}
		SELECT 
		f.departure_city,
		f.arrival_city,
		max(tf.amount) AS max_price,
		min(tf.amount) AS min_price
		FROM flights_v f
		JOIN ticket_flights tf ON f.flight_id = tf.flight_id
		GROUP BY f.departure_city, f.arrival_city
		ORDER BY f.departure_city, f.arrival_city;
	\end{lstlisting}
	
	\section{Выводы}
	В ходе выполнения лабораторной работы были освоены механизмы составления сложных SQL-запросов SELECT в СУБД PostgreSQL, научились применять различные типы соединений (JOIN), агрегатные функции, условные конструкции, работу с массивами и представление данных.
	
\end{document}